\chapter{Menata Data}\label{ch:g7:stats}

Sebelum menghitung apa pun, data harus dikumpulkan, dibilang dan
ditampilkan. Bab ini membahas cara mengubah daftar pengamatan yang
berantakan menjadi tabel atau gambar yang jelas: \omterm{def:g7:stats:frequency}{frekuensi}, \omterm{def:g7:stats:frequency}{frekuensi
relatif}, data yang dikelompokkan, diagram batang dan diagram
lingkaran. Merangkum data dengan satu bilangan (rata-rata, median)
menjadi bahasan \cref{ch:g9:statproba}.

\section{Frekuensi dan frekuensi relatif}

\begin{definition}[Frekuensi dan frekuensi relatif]\label{def:g7:stats:frequency}
Untuk setiap nilai yang mungkin pada sebuah survei,
\emph{frekuensinya}\index{frekuensi} adalah berapa kali nilai itu
muncul; \emph{frekuensi relatifnya}\index{frekuensi relatif} adalah
frekuensinya dibagi banyaknya seluruh pengamatan:
\[
\text{frekuensi relatif} = \frac{\text{frekuensi}}{\text{jumlah seluruhnya}} .
\]
Frekuensi relatif dapat ditulis sebagai \omterm{def:g6:fractions:def}{pecahan}, \omterm{def:g6:decimals:places}{bilangan desimal} atau
persen, dan \omterm{def:g1:addition:def}{jumlahnya} selalu $1$ (yaitu $100\,\%$).
\end{definition}

\begin{example}\label{ex:g7:stats:frequency}
Olahraga kesukaan $25$ murid sebuah kelas, dari daftar mentahnya
sampai menjadi tabel:
\begin{center}
\renewcommand{\arraystretch}{1.3}
\begin{tabular}{l|cccc|c}
olahraga & sepak bola & renang & tenis & judo & \omterm{def:g1:addition:def}{jumlah} \\
\hline
\omterm{def:g7:stats:frequency}{frekuensi} & $10$ & $6$ & $5$ & $4$ & $25$ \\
\omterm{def:g7:stats:frequency}{frekuensi relatif} & $0.40$ & $0.24$ & $0.20$ & $0.16$ & $1$ \\
\end{tabular}
\end{center}
Untuk sepak bola: $\frac{10}{25} = \frac{40}{100} = 0.40 = 40\,\%$.
Baris \omterm{def:g7:stats:frequency}{frekuensi relatifnya} harus berjumlah $1$ --- pemeriksaan tetap yang gratis.
\end{example}

\begin{remark}[Mengapa frekuensi relatif?]
\omterm{def:g7:stats:frequency}{Frekuensi} tidak dapat dibandingkan antara kelompok yang besarnya
berbeda: $10$ penggemar sepak bola di kelas berisi $25$ orang (yaitu
$40\,\%$) menunjukkan semangat yang lebih besar daripada $12$ orang di
satu angkatan berisi $120$ ($10\,\%$). \omterm{def:g7:stats:frequency}{Frekuensi relatif} menaruh
semuanya pada skala yang sama.
\end{remark}

\section{Mengelompokkan data ke dalam kelas}

\begin{example}[Kelas]\label{ex:g7:stats:classes}
Tinggi badan (dalam cm) $20$ murid berkisar dari $148$ sampai $172$.
Menyebutkan setiap tinggi secara terpisah akan memberi tabel dengan
kolom yang hampir sebanyak muridnya; sebagai gantinya, kelompokkanlah
ke dalam \emph{kelas}:
\begin{center}
\renewcommand{\arraystretch}{1.3}
\begin{tabular}{l|ccc|c}
tinggi (cm) & $\intco{145}{155}$ & $\intco{155}{165}$ &
$\intco{165}{175}$ & \omterm{def:g1:addition:def}{jumlah} \\
\hline
\omterm{def:g7:stats:frequency}{frekuensi} & $5$ & $9$ & $6$ & $20$ \\
\omterm{def:g7:stats:frequency}{frekuensi relatif} & $0.25$ & $0.45$ & $0.30$ & $1$ \\
\end{tabular}
\end{center}
Setiap tinggi terbilang di tepat satu kelas ($155$ masuk ke
$\intco{155}{165}$, bukan ke keduanya --- itulah sebabnya selangnya
setengah terbuka). Mengelompokkan kehilangan kerincian tetapi
memperoleh keterbacaan.
\end{example}

\section{Diagram}

\begin{method}[Memilih dan menggambar diagram]\label{met:g7:stats:charts}
\begin{enumerate}
  \item \emph{Diagram batang}: satu batang untuk tiap nilai atau
        kelas, tingginya \omterm{def:g6:propdata:def}{sebanding} dengan \omterm{def:g7:stats:frequency}{frekuensinya} (atau dengan
        \omterm{def:g7:stats:frequency}{frekuensi relatifnya}). Bagus untuk membandingkan nilai.
  \item \emph{Diagram lingkaran}: satu juring untuk tiap nilai,
        sudutnya \omterm{def:g6:propdata:def}{sebanding} dengan \omterm{def:g7:stats:frequency}{frekuensi relatifnya} --- seluruh
        lingkarannya ($360^\circ$) mewakili jumlah seluruhnya, jadi
        \[
        \text{sudut} = \text{frekuensi relatif} \times 360^\circ .
        \]
        Bagus untuk memperlihatkan bagian dari suatu keseluruhan.
  \item Pada kedua keadaan: berilah judul, nama, dan skala atau persen
        yang ditulis pada diagramnya.
\end{enumerate}
\end{method}

\begin{example}\label{ex:g7:stats:pie}
Sudut untuk olahraga pada \cref{ex:g7:stats:frequency}: sepak bola
$0.40 \times 360 = 144^\circ$; renang $0.24 \times 360 =
86.4^\circ$; tenis $72^\circ$; judo $57.6^\circ$. Periksa:
$144 + 86.4 + 72 + 57.6 = 360$.
\end{example}

\begin{omfigure}
\begin{tikzpicture}
\begin{axis}[
  width=6.6cm, height=5.2cm,
  ybar, bar width=14pt,
  symbolic x coords={sepakbola,renang,tenis,judo},
  xtick=data,
  x tick label style={font=\footnotesize, rotate=25, anchor=east},
  ymin=0, ymax=12,
  ytick={2,4,6,8,10},
  ylabel={frekuensi},
  label style={font=\small},
  tick label style={font=\small}]
  \addplot[fill=omDef!30, draw=omDef] coordinates
    {(sepakbola,10) (renang,6) (tenis,5) (judo,4)};
\end{axis}
\end{tikzpicture}\hspace{0.8cm}%
\begin{tikzpicture}[scale=1.05]
  \fill[omDef!45] (0,0) -- (0:1.5) arc (0:144:1.5) -- cycle;
  \fill[omThm!35] (0,0) -- (144:1.5) arc (144:230.4:1.5) -- cycle;
  \fill[omMeth!35] (0,0) -- (230.4:1.5) arc (230.4:302.4:1.5) -- cycle;
  \fill[omProp!35] (0,0) -- (302.4:1.5) arc (302.4:360:1.5) -- cycle;
  \draw (0,0) circle (1.5);
  \foreach \a in {0,144,230.4,302.4} \draw (0,0) -- (\a:1.5);
  \node[font=\footnotesize] at (72:0.95) {bola};
  \node[font=\footnotesize] at (187:0.95) {renang};
  \node[font=\footnotesize] at (266:0.95) {tenis};
  \node[font=\footnotesize] at (331:0.95) {judo};
\end{tikzpicture}

{\small Data yang sama sebagai diagram batang (bandingkan tingginya)
dan sebagai diagram lingkaran (bandingkan bagian dari kelasnya). Sepak
bola mengambil $144^\circ$, yaitu $40\,\%$ lingkarannya.}
\end{omfigure}

\begin{example}[Membaca secara kritis]\label{ex:g7:stats:critical}
Diagram yang sumbu tegaknya bermula dari $9$ dan bukan dari $0$
membuat batang $10$ terlihat sepuluh kali lebih tinggi daripada batang
$9.1$ --- menyesatkan mata, walaupun bilangannya jujur. Naluri pertama
di depan diagram apa pun: periksalah di mana sumbunya bermula dan
berapa nilai satu skalanya (\cref{met:g6:propdata:read}).
\end{example}

\section{Latihan}

\begin{exercise}[$\star$]\label{exo:g7:stats:1}
Berikut nilai $20$ murid: $12$, $15$, $9$, $12$, $14$,
$9$, $12$, $15$, $11$, $12$, $9$, $14$, $12$, $11$, $15$, $9$, $12$,
$14$, $11$, $12$. Susunlah tabel \omterm{def:g7:stats:frequency}{frekuensinya} (satu kolom untuk tiap nilai).
\end{exercise}

\begin{exercise}[$\star$]\label{exo:g7:stats:2}
Lengkapilah tabel pada \cref{exo:g7:stats:1} dengan baris \omterm{def:g7:stats:frequency}{frekuensi
relatif} (sebagai \omterm{def:g6:fractions:def}{pecahan} dari $20$, lalu sebagai persen). Periksalah \omterm{def:g1:addition:def}{jumlahnya}.
\end{exercise}

\begin{exercise}[$\star$]\label{exo:g7:stats:3}
Pada survei terhadap $50$ keluarga, $30$ keluarga punya satu mobil.
Berapa \omterm{def:g7:stats:frequency}{frekuensi relatif} ``satu mobil'', sebagai \omterm{def:g6:decimals:places}{bilangan desimal} dan
sebagai persen? Pada survei yang lain, $45$ keluarga dari $90$ punya
satu mobil: survei mana yang menunjukkan proporsi yang lebih tinggi?
\end{exercise}

\begin{exercise}[$\star$]\label{exo:g7:stats:4}
\omterm{def:g2:measure:units}{Massa} (kg) $15$ ekor anjing: $8$, $23$, $31$, $12$, $9$, $27$, $35$,
$14$, $22$, $18$, $29$, $11$, $25$, $33$, $16$. Kelompokkan ke dalam
kelas $\intco{5}{15}$, $\intco{15}{25}$, $\intco{25}{35}$,
$\intco{35}{45}$ lalu sebutkan \omterm{def:g7:stats:frequency}{frekuensinya}.
\end{exercise}

\begin{exercise}[$\star$]\label{exo:g7:stats:5}
Gambarlah diagram batang dari tabel ini:
\begin{center}
\renewcommand{\arraystretch}{1.2}
\begin{tabular}{l|cccc}
hewan peliharaan & $0$ & $1$ & $2$ & $3$ \\
\hline
keluarga & $8$ & $12$ & $6$ & $4$ \\
\end{tabular}
\end{center}
\end{exercise}

\begin{exercise}[$\star$]\label{exo:g7:stats:6}
Untuk tabel pada \cref{exo:g7:stats:5}, hitunglah \omterm{def:g7:stats:frequency}{frekuensi relatif}
dan sudut diagram lingkaran tiap nilainya (periksa: \omterm{def:g1:addition:def}{jumlahnya} $360^\circ$).
\end{exercise}

\begin{exercise}[$\star$]\label{exo:g7:stats:7}
Sebuah diagram lingkaran tentang musim kesukaan menunjukkan: musim
panas $162^\circ$, musim semi $90^\circ$, musim gugur $54^\circ$,
musim dingin $54^\circ$. Berapa persen yang memilih tiap musim? Kalau
$60$ orang menjawab, berapa yang memilih musim panas?
\end{exercise}

\begin{exercise}[$\star\star$]\label{exo:g7:stats:8}
Di kelas A, $12$ murid dari $30$ berjalan kaki ke sekolah; di kelas B,
$14$ dari $40$. Kelas mana yang proporsi pejalan kakinya lebih tinggi?
Berilah alasan dengan \omterm{def:g7:stats:frequency}{frekuensi relatif}, bukan dengan \omterm{def:g7:stats:frequency}{frekuensi}.
\end{exercise}

\begin{exercise}[$\star\star$]\label{exo:g7:stats:9}
Sebuah majalah mencetak diagram batang penjualan bulanannya:
$9\,800$, lalu $10\,000$, lalu $10\,100$, dengan sumbu tegaknya
bermula dari $9\,700$. Jelaskan apa yang dilihat pembacanya, dan apa
yang akan diperlihatkan sumbu yang jujur yang bermula dari $0$.
\end{exercise}

\begin{exercise}[$\star\star$]\label{exo:g7:stats:10}
Sebuah tabel \omterm{def:g7:stats:frequency}{frekuensi relatif} punya tiga nilai dengan \omterm{def:g7:stats:frequency}{frekuensi
relatif} $0.35$, $0.4$ dan $f$. Carilah $f$. Kalau \omterm{def:g1:addition:def}{jumlah} \omterm{def:g7:stats:frequency}{frekuensinya}
$80$, sebutkan ketiga \omterm{def:g7:stats:frequency}{frekuensinya}.
\end{exercise}

\begin{exercise}[$\star\star\star$]\label{exo:g7:stats:11}
Di sebuah sekolah, $55\,\%$ muridnya perempuan. Di antara yang
perempuan, $40\,\%$ makan di kantin; di antara yang laki-laki,
$60\,\%$ makan di kantin. Dari $400$ murid, berapa yang makan di
kantin? (Hitunglah besar keempat kelompoknya selangkah demi
selangkah.) Berapa persennya secara keseluruhan?
\end{exercise}

\section{Soal: Bagaimana data berbohong --- dan cara menangkapnya}

\begin{problem}[{Soal akhir pekan --- frekuensi melawan frekuensi
relatif, jajak pendapat yang menipu satu negara, dan paradoks dua
rumah sakit}]
\label{pb:g7:stats:1}
Bilangan tidak berbohong, tetapi bilangan dapat dibuat menyesatkan:
\omterm{def:g7:stats:frequency}{frekuensi} yang besar dapat menyembunyikan \omterm{def:g7:stats:frequency}{frekuensi relatif} yang
kecil, survei yang raksasa dapat tidak berguna, dan --- yang paling
aneh --- sebuah rumah sakit dapat mengalahkan saingannya pada
\emph{setiap} golongan pasien dan tetap kalah pada angka
keseluruhannya. Ketiga perangkap itu dilepaskan pada soal ini, hanya
berbekal \omterm{def:g7:stats:frequency}{frekuensi} dan \omterm{def:g7:stats:frequency}{frekuensi relatif} bab ini
(\cref{def:g7:stats:frequency}).

\medskip
\noindent\textbf{Bagian I --- \omterm{def:g7:stats:frequency}{Frekuensi} bukan \omterm{def:g7:stats:frequency}{frekuensi relatif}.}
\begin{enumerate}
  \item Sekolah A mendaur ulang $120$ dari $400$ kotak jusnya; sekolah
        B mendaur ulang $90$ dari $250$ kotaknya. Sekolah mana yang
        \emph{\omterm{def:g7:stats:frequency}{frekuensinya}} lebih besar? Yang \emph{\omterm{def:g7:stats:frequency}{frekuensi
        relatifnya}} lebih besar? Sekolah mana yang layak menerima
        hadiah daur ulang?
  \item Sebuah jadwal baru disurvei: $18$ dari $30$ gurunya menyukai
        jadwal itu, dan $210$ dari $600$ muridnya menyukainya.
        Hitunglah kedua \omterm{def:g7:stats:frequency}{frekuensi relatifnya}, lalu \omterm{def:g7:stats:frequency}{frekuensi relatif}
        pada seluruh $630$ orang itu. Mengapa angka gabungannya
        terletak begitu dekat dengan angka para muridnya?
  \item Sebuah diagram lingkaran yang diterbitkan menunjukkan juring
        berlabel $45\,\%$, $30\,\%$, $20\,\%$ dan $10\,\%$. Tanpa
        keterangan lain, bagaimana kamu tahu ada kekeliruan?
  \item Sebuah tabel \omterm{def:g7:stats:frequency}{frekuensi relatif} atas $40$ pengamatan terhapus
        separuhnya: \omterm{def:g7:stats:frequency}{frekuensinya} $14$, $10$, $?$, $?$ dan \omterm{def:g7:stats:frequency}{frekuensi
        relatifnya} $0.35$, $0.25$, $0.15$, $?$. Bangunlah kembali
        entri yang hilang.
  \item Sebuah survei berakhir dengan \omterm{def:g7:stats:frequency}{frekuensi relatif} $0.40$, $0.35$
        dan $0.25$. Hitunglah ketiga sudut diagram lingkarannya
        (\cref{met:g7:stats:charts}) lalu periksalah bahwa ketiganya
        menutup lingkarannya.
\end{enumerate}

\medskip
\noindent\textbf{Bagian II --- Jajak pendapat yang menipu satu negara.}
Pada tahun 1936, sebuah majalah Amerika mengirimkan sepuluh juta surat
suara ke alamat yang diambil dari buku telepon dan daftar pendaftaran
mobil, menerima $2.4$ juta jawaban, lalu meramalkan kemenangan telak
untuk calon bernama Landon. Seorang ahli statistik muda, George
Gallup, bertanya hanya kepada kira-kira lima puluh ribu orang --- yang
dipilih agar menyerupai seluruh penduduknya --- lalu meramalkan
sebaliknya. Roosevelt menang dengan telak.
\begin{enumerate}[resume]
  \item Berapa \omterm{def:g7:stats:frequency}{frekuensi relatif} surat suara yang kembali?
  \item Pada tahun 1936 telepon dan mobil masih barang mewah.
        Jelaskan dalam satu atau dua kalimat mengapa sampel majalah
        itu, walaupun raksasa, sudah pasti gagal --- dan pelajaran
        pertanyaan 1 yang mana yang diulanginya pada skala nasional.
  \item Sebuah model mainan. Sebuah kota punya $1\,000$ pemilih kaya,
        yang $70\,\%$ di antaranya mendukung L, dan $9\,000$ pemilih
        sederhana, yang $30\,\%$ di antaranya mendukung L. Hitunglah
        dukungan L yang sebenarnya di kota itu. Sebuah ``jajak
        pendapat buku telepon'' menjangkau $500$ pemilih kaya dan
        $500$ pemilih sederhana: dukungan berapa yang diramalkannya?
  \item Perangkap yang diam: sebuah perusahaan punya $200$ pelanggan
        yang tidak puas dan $800$ pelanggan yang puas. Ia mensurvei
        semuanya; $40\,\%$ yang tidak puas menjawab, tetapi hanya
        $10\,\%$ yang puas menjawab. Hitunglah banyaknya jawaban tiap
        jenisnya, dan \omterm{def:g7:stats:frequency}{frekuensi relatif} ketidakpuasan \emph{di antara
        jawabannya}. Bandingkan dengan \omterm{def:g7:stats:frequency}{frekuensi relatif} yang sebenarnya.
  \item Perbaikilah jajak pendapat pada pertanyaan 8: pertahankan
        $1\,000$ wawancara, tetapi bagikan menurut proporsi kota yang
        sebenarnya. Apa yang diramalkan jajak pendapat yang sudah diperbaiki itu?
\end{enumerate}

\medskip
\noindent\textbf{Bagian III --- Paradoks dua rumah sakit.}
Dua rumah sakit menerbitkan \omterm{def:g1:addition:def}{jumlah} pasien yang sembuh, dipilah
menurut tingkat keparahannya:
\begin{center}
\renewcommand{\arraystretch}{1.3}
\begin{tabular}{l|cc}
 & kasus ringan & kasus berat \\
\hline
rumah sakit A & $90$ sembuh dari $100$ & $280$ sembuh dari $400$ \\
rumah sakit B & $340$ sembuh dari $400$ & $65$ sembuh dari $100$ \\
\end{tabular}
\end{center}
\begin{enumerate}[resume]
  \item Hitunglah \omterm{def:g7:stats:frequency}{frekuensi relatif} kesembuhan rumah sakit A untuk
        kasus ringan, untuk kasus berat, dan secara keseluruhan
        (seluruh $500$ pasiennya).
  \item Tiga perhitungan yang sama untuk rumah sakit B. Rumah sakit
        mana yang menang pada kasus ringan? Pada kasus berat?
  \item Bandingkan kedua \omterm{def:g7:stats:frequency}{frekuensi relatif} keseluruhannya. Nyatakan
        dengan gamblang hal aneh yang baru saja terjadi.
  \item Jelaskan triknya: bandingkan \emph{campuran} pasien kedua
        rumah sakit itu, lalu katakan mengapa \omterm{def:g7:stats:frequency}{frekuensi relatif}
        keseluruhannya dapat mengkhianati kedua \omterm{def:g7:stats:frequency}{frekuensi relatif}
        golongannya. Kalau kamu terkena kasus yang berat, rumah sakit
        mana yang sebaiknya kamu pilih?
  \item Sebuah judul berita berbunyi: ``Rumah sakit B punya angka
        kesembuhan yang lebih baik: $81\,\%$ melawan $74\,\%$.''
        Tulislah surat dua kalimat kepada redaksi yang sudah diajarkan
        seluruh soal ini kepadamu --- satu kalimat tentang apa yang
        diabaikan judul itu, satu kalimat tentang pesan moral umumnya
        (bandingkan yang sejenis dengan yang sejenis, dan perhatikan
        bobotnya).

\end{enumerate}
\end{problem}
